\begin{figure*}[t]
\centering
\includegraphics[width=\textwidth]{Figure_B1_fixed_task0_gradient_near_origin.pdf}
\caption{Fixed Task-0 positive-gradient response on the formal ImageNet-A seeds 0, 1, and 2 (ViT-B/16, rank-10 Q/V LoRA). At each method/seed own Task-0 checkpoint, compute once the gradient g of mean true-label cross-entropy over all 1,259 Task-0 training examples with respect to the 24 effective Q/V weight matrices, then set u=+g/||g||. Evaluate theta(r)=theta(0)+r*u with this single direction held fixed for the whole radius scan; the old 20-class head and the remaining model state are fixed. Each method uses its own gradient direction; directions are not shared between EWCLoRA and FS-EWCLoRA. Left: signed Task-0 training CE difference D\_0(r)=L\_0(theta(r))-L\_0(theta(0)), retaining any negative values without clipping. Middle: Task-0 training accuracy (\%). Right: raw empirical Fisher trace on the unchanged fixed 256-example training subset, computed from per-example true-label CE gradients in the same effective Q/V coordinates. Fisher has no extra normalization, floor, regularizer, or multiplication by the EWC coefficient. The full scan contains 93 measured radii: 0–20 in steps of 0.25, 25, 30, 35, 40, 45, 50, and additional near-origin radii 0.001, 0.005, 0.01, 0.025, 0.05, 0.1. This companion shows only the measured 0–2 window to reveal the initial response. Lines and shading are the arithmetic mean and one sample standard deviation (ddof=1) across the three training seeds. Each seed/method supplies one direction; directions are not additional independent replicates. Connecting segments are visual guides, with no interpolation treated as a measurement. The CSV also exports max(D\_0(r),0), alongside signed training/test CE changes and accuracy, but the displayed loss panel uses signed D\_0. Gradient unit norm, saved direction identity, full Task-0 sample membership, formal origin recovery, and an initial finite-difference ascent witness are checked before publishing. This is a fixed direction in effective weights, not an actual LoRA optimizer update path, iterative gradient descent, a Task-1 update, or an experiment measuring future-task forgetting.}
\label{fig:fixed-task0-gradient-near-origin}
\end{figure*}
